%% ====================================================================
\section{Hipótesis de partida}
\label{sec:hipotesis}
%% ====================================================================

\textbf{Pregunta de investigación.}
¿Puede una arquitectura distribuida basada en dinámicas de Smith formar coaliciones
multi-AGV factibles en un almacén con comunicación local, tareas heterogéneas y demanda
temporal, y qué mecanismo es necesario cuando la comunicación se degrada hasta romper el
supuesto de conectividad suficiente?

La respuesta se estructura en hipótesis mecanísticas. Cada hipótesis se contrasta con
evidencia reproducible y los resultados negativos se reportan con el mismo peso que los
positivos.

\textbf{H1: Smith reproduce water-filling en equilibrio homogéneo.}
En el motor estático, cuando las tareas comparten estructura y no hay restricciones
temporales, Smith asigna la población según la regla de \emph{water-filling}. El gate
reportado pasa con pendiente $1.0000$ y $R^2=1$ en el protocolo mínimo.

\textbf{H2: el clearing entero reduce desperdicio físico.}
La asignación fluida puede repartir masa sobre tareas que no alcanzan quorum entero. El
\emph{clearing} entero debe reducir distancia desperdiciada y coaliciones parciales. En
el protocolo mínimo reduce el desperdicio alrededor de un $32\%$.

\textbf{H3-A: el precio no es una mejora estática.}
En equilibrio estático, añadir precio a Smith puede restar valor porque perturba un
reparto que ya era óptimo. Este resultado negativo queda documentado y no se usa como
fallo de la teoría.

\textbf{H3-B: el precio solo tiene hipótesis válida en el régimen temporal.}
Con llegadas, \emph{deadlines} y déficit persistente, el precio puede rescatar tareas
antes de expirar. Esta es la formulación correcta de la hipótesis de precios y se valida
en el motor temporal aislado.

\textbf{H4: Smith original falla bajo R3.}
Cuando $R_{\mathrm{com}}$ cae a R3, los robots pueden conocer y perseguir tareas sin
cerrar quorums físicos suficientes. La dinámica se paraliza no porque sea distribuida,
sino porque su equilibrio distribuido necesita conectividad local suficiente.

\textbf{H5: Smith-QR rescata parcialmente el colapso de Smith en R3.}
Smith-QR debe mejorar a Smith original bajo R3 mediante memoria de descriptores,
compromiso en ruta, guardia de \emph{clearing} y proyección de quorum local. En la
validación controlada de 20 semillas pasa de $0.1038$ a $0.2893$ de captura media frente
a Smith; el delta pareado es $+0.1855$ con IC95 $[0.1434,0.2276]$.

\textbf{H6: Smith-QR solo se declara dominante en métricas con IC95 positivo.}
La campaña R3 de 20 semillas sí muestra captura superior a greedy y al proxy MARL inicial, pero no
convierte a Smith-QR en dominante universal. La distancia desperdiciada y el coste
operativo deben reportarse por separado; si una métrica no tiene delta pareado positivo
con IC95 disjunto de cero, se interpreta como paridad o coste de robustez.

\textbf{H7: CoppeliaSim valida plausibilidad robótica, no sustituye el benchmark.}
Las escenas industriales con Pioneer~3-DX, humanos, sensores y racks verifican que la
historia robótica es plausible. La decisión científica principal sigue viniendo del
benchmark reproducible.

\textbf{H8: MARL-CTDE debe tratarse como baseline aprendido, no como sustituto teórico.}
Una política MARL con entrenamiento centralizado y ejecución descentralizada puede
aprender pesos compartidos para puntuar pares robot--carga a partir de observaciones
locales. La hipótesis defendible no es que MARL domine a Smith-QR, sino que el comparador
de aprendizaje sea real: debe mostrar ganancia durante entrenamiento, evaluarse con
pesos congelados y reportar cuándo mejora al proxy y cuándo falla frente a métodos
analíticos.

\textbf{H9: el actor neuronal CTDE aporta comparación aprendida adicional, no una
sustitución automática.}
Un actor neuronal compartido puede ampliar la política lineal mediante una capa no lineal
residual entrenada con retornos episódicos de la flota. La hipótesis no exige que el
actor neuronal venza a Smith-QR; exige que aprenda bajo el mismo protocolo, que su
evaluación congelada mejore al proxy en algún régimen y que cualquier mejora frente al
actor lineal se declare sólo con métricas e intervalos de confianza explícitos.

\textbf{H10: la densidad predictiva reduce congestión sin degradar productividad.}
El peaje de congestión inducido por $\rho(x,t+\tau)$ debe desviar parte del flujo antes
de que el cuello de botella se bloquee. La hipótesis se acepta si Smith-QR con densidad
predictiva reduce pasos congestionados y densidad futura máxima frente a Smith-QR base,
manteniendo captura de recompensa no inferior dentro de un margen predefinido.

\textbf{H11: la capacidad wrench mejora el tramo dinámico crítico.}
Una coalición evaluada por cardinalidad puede no producir el wrench requerido. La
hipótesis se acepta si la política con capacidad wrench reduce residual físico frente a
cardinalidad simple y mantiene cotas de torque, salto Hamiltoniano y batería durante el
transporte de una carga rectangular.

\textbf{H12: CoppeliaSim aporta cobertura visual reproducible de los escenarios de
validación.}
La hipótesis no declara ejecución hardware ni sustituye la estadística principal. Pasa
si cada escenario necesario para el TFM queda exportado como escena CoppeliaSim con
cámaras de auditoría, datos de playback y artefactos visuales replicables, o con un
render sintético trazable cuando no se ejecuta CoppeliaSim en modo headless.

\subsection{Criterios cuantitativos de contraste}
\label{subsec:criterios}

Cada hipótesis se interpreta con una hipótesis nula conservadora $H_0$ y una
alternativa $H_1$. Cuando existe una demostración cerrada, $H_0$ se refuta por argumento
matemático y el gate numérico sólo audita la implementación algebraica. Cuando la
afirmación depende de repeticiones, semillas o comparación entre tratamientos, la
decisión se toma con contrastes unilaterales o intervalos de confianza pareados al nivel
$\alpha=0{,}05$. La validación CoppeliaSim se reporta como evidencia robótica visual con
datos exportados de escena; no se confunde con inferencia poblacional ni con hardware
real.

\begin{description}
  \item[H1:] pendiente de ajuste lineal entre staffing teórico y observado en
  $[0.99,1.01]$ y $R^2\geq0.99$.
  \item[H2:] delta pareado de desperdicio menor que cero al activar
  \emph{clearing} entero.
  \item[H3-A:] se acepta como negativo si precio estático no mejora a Smith sin precio.
  \item[H3-B:] precio temporal pasa si mejora captura antes de \emph{deadline} con
  IC95 positivo en delta pareado.
  \item[H5:] Smith-QR pasa si en R3 mejora a Smith original con IC95 positivo.
  \item[H6:] dominancia frente a greedy/MARL solo se afirma si el delta pareado tiene
  IC95 positivo; si no, se documenta como límite operativo.
  \item[H8:] MARL-CTDE pasa como baseline aprendido si aumenta el objetivo episódico
  durante entrenamiento y su evaluación congelada mejora al proxy en al menos un régimen
  con IC95 positivo; cualquier no-dominancia frente a Smith-QR se reporta como límite.
  \item[H9:] el actor neuronal CTDE pasa como baseline neuronal si aumenta el objetivo
  episódico, mejora al proxy en escasez y reduce al menos un coste operativo frente al
  actor lineal con IC95 favorable; si no domina a Smith-QR, se reporta como límite.
  \item[H10:] densidad predictiva pasa si reduce pasos congestionados y densidad futura
  máxima con IC95 favorable, sin degradar captura más de un margen de no inferioridad.
  \item[H11:] motor integrado pasa si reduce residual wrench y mantiene torque,
  Hamiltoniano y batería dentro de cotas operativas.
  \item[H12:] CoppeliaSim pasa como cobertura visual si todos los escenarios exportan
  escena, cámaras, playback, PNG y MP4 o fallback sintético declarado.
  \item[CoppeliaSim:] una escena pasa el smoke si abre, contiene robots Pioneer o
  fallback equivalente, sensores, humanos, racks y produce métricas de colisión,
  distancia mínima, saturación y eventos de quorum/entrega.
\end{description}
